Jupiter is the biggest planet in our Solar System. Its mass and radius are
$M_J = 1.90 \times 10^{27}\,\mathrm{kg}$ and $R_J = 6.99 \times 10^{7}\,\mathrm{m}$,
respectively. Herein, we investigate Jupiter's internal structure using a
simple model. As shown in Figure T11.1, this gas giant is assumed to contain
two distinct regions: a dense core and a fluid envelope. While the dense core
is entirely made of rock, the fluid envelope is wholly composed of hydrogen.
The mass densities of rock $(\rho_Z)$ and hydrogen $(\rho_H)$ are given by
\[ \rho_Z = \mathrm{const}, \quad \rho_H = \sqrt{\frac{P}{K}} \]
where $P$ is the hydrostatic pressure and $K$ is the polytropic coefficient.
The space outside Jupiter is considered an absolute vacuum. The effects of
self-rotation are neglected to achieve perfect spherical symmetry. These
approximations enable us to study Jupiter's structural properties via a single
spatial variable: the distance $r$ from the center $O$.

\ioaafig{2026_TH_T11-f1.pdf}

\medskip\noindent\textbf{Part 1 (23 points):}

Assume that the dense core is not eroded by the fluid envelope. In this case,
$P$ decreases from $P_O = 7.13 \times 10^{12}\,\mathrm{Pa}$ at the center
$(r = 0)$ to $P_C = 3.98 \times 10^{12}\,\mathrm{Pa}$ at the core-envelope
boundary $(r = R_C = 1.05 \times 10^{7}\,\mathrm{m})$. The polytropic
coefficient of pure hydrogen is
$K = 2.12 \times 10^{5}\,\mathrm{m^5\,kg^{-1}\,s^{-2}}$.

\begin{description}
\item[(a)] (6 points) Find the value of $\rho_Z$.
\item[(b)] (6 points) Since the fluid envelope is equivalent to a simple
  polytrope $(P \propto \rho_H^2)$, its density depends on the distance as
  \[ \rho_H = \frac{A}{r}\cos\left(\frac{r}{R} + \varphi\right) \]
  where $R$ is the characteristic radius, $A$ is the amplitude factor, and
  $\varphi$ is the phase constant. Given that
  \[ R = 2.25 \times 10^{7}\,\mathrm{m}, \quad A > 0, \quad
     -\pi \le \varphi \le \pi \]
  Determine $A$ and $\varphi$.
\item[(c)] (8 points) Plot the density and pressure versus the distance for
  the whole of Jupiter's interior. Each graph must consist of at least six
  distinct points.
\item[(d)] (3 points) Calculate the mass of the dense core $M_Z$ and the mass
  of the fluid envelope $M_H$.
\end{description}

\medskip\noindent\textbf{Part 2 (27 points):}

Recent studies show that giant planets may undergo core erosion due to the
high solubility of rock in hydrogen. For Jupiter, this process helps
redistribute heavy and light materials throughout its interior. When the dense
core disappears, Jupiter becomes an ideal mixture of rock and hydrogen at
equilibrium. Neglect all possible chemical reactions. For a given position
inside Jupiter, the mass fraction of rock $Z$ is defined by
\[ Z = \frac{\mathrm{d}m_Z}{\mathrm{d}m_Z + \mathrm{d}m_H} \]
where $\mathrm{d}m_Z$ and $\mathrm{d}m_H$ are the differential masses of rock
and hydrogen, respectively.

\begin{description}
\item[(a)] (6 points) Let $\rho$ be the mass density of the mixture at the
  distance $r$. Prove that $\rho$ is related to $P$ via
  \[ \rho = \frac{\rho_Z\sqrt{P}}{Z\sqrt{P} + (1 - Z)\rho_Z\sqrt{K}} \]
\item[(b)] (7 points) Figure T11.2 shows how $Z$ depends on $r$. Based on the
  $Z - r$ relation, we can divide Jupiter's interior into three different
  regions: the dilute core $(Z = Z_1 = 7.28 \times 10^{-2},\ 0 \le r \le R_1)$,
  the transition layer ($Z$ decreases linearly with increasing $r$,
  $R_1 \le r \le R_2$), and the hydrogen-dominated envelope
  $(Z = Z_2 = 1.53 \times 10^{-2},\ R_2 \le r \le R_J)$. In the transition
  layer, $\rho$ is well approximated as a linear function of $r$, and $P$
  drops from $P_1 = 2.25 \times 10^{12}\,\mathrm{Pa}$ at
  $R_1 = 2.87 \times 10^{7}\,\mathrm{m}$ to
  $P_2 = 8.43 \times 10^{11}\,\mathrm{Pa}$ at
  $R_2 = 4.40 \times 10^{7}\,\mathrm{m}$. The polytropic coefficient of
  hydrogen in the mixture is
  $K = 2.20 \times 10^{5}\,\mathrm{m^5\,kg^{-1}\,s^{-2}}$. Determine the mass
  of the transition layer $M_{12}$.

  \ioaafig{2026_TH_T11-f2.pdf}

\item[(c)] (14 points) Calculate the mass of the dilute core $M_1$ and the
  mass of the hydrogen-dominated envelope $M_2$.
\end{description}
